\section{Manual de Uso}

Esta secci\'on detalla los pasos necesarios para integrar la librer\'ia
\texttt{orange-jumpsuit} en proyectos de C y Python, as\'i como el flujo de
trabajo recomendado para realizar b\'usquedas vectoriales.

\subsection{Integraci\'on en C}

La librería est\'a dise\~nada como un archivo de cabecera \'unico
(\textit{header-only}), lo que facilita su portabilidad entre 
proyectos.

\subsubsection{Requisitos y Dependencias}
Para compilar un proyecto que utilice \texttt{jumpsuit.h}, se requiere:
\begin{itemize}
    \item Un compilador compatible con C11 o superior.
    \item La librería \textbf{OpenBLAS} (u otra implementaci\'on de BLAS) 
      instalada en el sistema para las operaciones matriciales aceleradas.
    \item En sistemas Linux, el script de construcci\'on de pruebas 
      \texttt{build\_tests.sh} puede servir como referencia para la 
      configuraci\'on del enlazador.
\end{itemize}

\subsubsection{Inclusi\'on en el Proyecto}
Para utilizar la librer\'ia, simplemente incluya el archivo en su c\'odigo
fuente. 
En un archivo de C, debe definir \verb|JUMPSUIT\_IMPLEMENTATION|
antes de la inclusi\'on para generar la implementaci\'on de las funciones:

\begin{lstlisting}[language=C]
#define JUMPSUIT_IMPLEMENTATION
#include "jumpsuit.h"
\end{lstlisting}

\subsection{Flujo de Trabajo en C}

El uso del \'indice sigue un ciclo de vida de cuatro etapas principales:
inicialización, entrenamiento, adición de datos y búsqueda.

\begin{enumerate}
    \item \textbf{Inicialización}: 
      Se definen las dimensiones del vector, el número de subespacios ($m$)
      y la cantidad de bits por centroide ($k^*$).
    \item \textbf{Entrenamiento}: 
      Se provee un conjunto de vectores para calcular los centroides del
      \verb|codebook|. 
      Se recomienda que este conjunto sea representativo de los datos 
      finales.
    \item \textbf{Adici\'on}: 
      Se cargan los vectores que formar\'an parte de la base de datos de 
      b\'usqueda. En este punto, los vectores se comprimen a sus \'indices
      correspondientes.
    \item \textbf{B\'usqueda}: 
      Se realiza la consulta para obtener los $k$ vecinos m\'as cercanos.
\end{enumerate}

\subsection{Integraci\'on en Python}

Se provee un wrapper que puede instalarse directamente vía \texttt{pip} en
sistemas Linux y Windows:

\begin{lstlisting}[language=bash]
pip install orange-jumpsuit
\end{lstlisting}

\subsubsection{Ejemplo de Uso en Python}
El wrapper de Python abstrae la gesti\'on de memoria y permite interactuar con
la librer\'ia utilizando arreglos de NumPy:

\begin{lstlisting}[language=python]
from orange_jumpsuit import IndexPQ

# 1. Configuracion del indice
index = IndexPQ(dimension=128, n_subvectors=8, n_bits_per_value=8)

# 2. Entrenamiento y carga de datos
index.train(data_entrenamiento)
index.add(data_base)

# 3. Busqueda de los 10 vecinos mas cercanos
distancias, indices = index.search(query_vector, n_neighbours=10)
\end{lstlisting}

El wrapper esta hecho mayoritariamente con IA. El link a la conversaci\'on:
\url{https://gemini.google.com/share/62c402d547e8}.

\newpage

Tambi\'en es importante aclarar que el compilado y publicado a PiPY se hace 
mediante \href{https://cibuildwheel.pypa.io/en/stable/}{CiBuildWheels} y Github 
Actions. Los documentos importantes son los siguientes: 

\dirtree{%
.1 orange-jumpsuit.
.2 .github.
.3 workflows. 
.4 publish.yml.  
.2 python.
.3 dummy.c. 
.3 MANIFEST.in. 
.3 pyproject.toml. 
.2 build\_python\_wheel.sh. 
.2 build\_python\_wheel.bat. 
}

Como se compila en maquinas virtuales en GitHub, es posible que no se encuentre
el \texttt{.whl} (programa compilado) correcto para la computadora del usuario
publicado. 
Dado esto, \texttt{pip} intenta compilar el programa con 
\texttt{gcc/cl.exe}, entonces deben estar instalados y accesibles en la computadora.

Como alternativa, se puede compilar con los scripts 
\texttt{build\_python\_wheel[.bat/sh]} y luego correr 
\begin{lstlisting}[language=bash]
pip install [camino a ./dist/.whl]
\end{lstlisting}
en el entorno de python.

De nuevo es importante aclarar que el proceso de publicaci\'on y compilado para 
python fue hecho con ayuda de IA. El link a la conversaci\'on: 
\url{https://gemini.google.com/share/4805b8db0b70}

\subsection{Pruebas y Verificaci\'on}

Para verificar la correcta instalaci\'on en entornos Linux, 
se pueden descargar los datos (\texttt{siftsmall.tar.gz} en 
\texttt{/datos/links.txt})
y compilar el ejecutable de pruebas incluido en el repositorio:

\begin{lstlisting}[language=bash]
./build_tests.sh
./build/tests/orange-jumpsuit
\end{lstlisting}

Este ejecutable realiza una prueba sobre las funciones expuestas por la 
librer\'ia.

\newpage

\subsection{Ejemplos}

El proyecto incluye dos \textit{Jupyter Notebooks} diseñados para demostrar 
la utilidad pr\'actica de la librer\'ia \texttt{orange\_jumpsuit} y evaluar
su rendimiento frente a soluciones est\'ndar de la industria como FAISS. 
A continuaci\'on se detalla el prop\'osito y funcionamiento de cada uno.

\subsection*{\texttt{SemanticSearch.ipynb}}

Este notebook es c\'odigo levemente modifciado de un ejemplo de busqueda 
sem\'antica con FAISS hecho por \textit{huggingface}:
\url{https://huggingface.co/learn/llm-course/chapter5/6} \\

\textbf{Prop\'osito:} 
Ilustrar un caso de uso real en Procesamiento de Lenguaje Natural (NLP),
demostrando c\'omo integrar la librer\'ia para implementar un motor de 
b\'usqueda sem\'antica de documentos.

\textbf{Funcionalidad del notebook}
\begin{itemize}
    \item \textbf{Carga de Datos:} 
      Utiliza la librer\'ia \texttt{datasets} de HuggingFace para descargar 
      y filtrar un corpus de texto real (ej. \textit{issues} y comentarios de
      GitHub).
    \item \textbf{Generaci\'on de Embeddings:} 
      Emplea \texttt{transformers} y \texttt{PyTorch} con el modelo
      pre-entrenado \texttt{multi-qa-mpnet-base-dot-v1} para convertir el
      texto en vectores densos (dimensi\'on 768).
    \item \textbf{Indexaci\'on y B\'usqueda:} 
      Construye \'indices de cuantizaci\'on de producto (PQ) usando tanto FAISS
      como \texttt{orange\_jumpsuit}.
    \item \textbf{Comparativa Pr\'actica:} 
      Ejecuta una consulta de ejemplo 
      (\textit{"How can I load a dataset offline?"}),
      mide los tiempos de ejecuci\'on de ambas herramientas y muestra
      visualmente que los resultados recuperados (Top-5) por
      \texttt{orange\_jumpsuit} son consistentes y relevantes.
\end{itemize}

\newpage 

\subsection*{\texttt{ComparingTest.ipynb}}

\textbf{Prop\'osito:} 
Realizar una prueba de rendimiento que compara la precisi\'on y velocidad de 
\texttt{orange\_jumpsuit} frente a FAISS.

\textbf{Funcionalidad del notebook}
\begin{itemize}
    \item \textbf{Dataset Est\'andar:}
      Carga el conjunto de datos \texttt{siftsmall}\footnote{
      Estos datos se encuentran en \url{http://corpus-texmex.irisa.fr/}}, 
      utilizado en la academia y la industria para evaluar estos algoritmos de
      b\'usqueda aproximada de vecinos m\'as cercanos.
    \item \textbf{Lectura de Archivos Binarios:} 
      Implementa funciones para leer los vectores flotantes (\texttt{.fvecs})
      y las respuestas correctas o \textit{groundtruth} (\texttt{.ivecs}).
    \item \textbf{Evaluaci\'on de Rendimiento (Tiempos):}
      Mide y compara estrictamente los tiempos de entrenamiento de los 
      centroides, la indexaci\'on de los vectores base y el tiempo de respuesta
      en las consultas (b\'usqueda).
    \item \textbf{Evaluaci\'on de Precisi\'on (M\'etricas):} 
      Analiza la calidad del \'indice calculando el \textbf{Recall@K} 
      (\texttt{Recall@1}, \texttt{Recall@10} y \texttt{Recall@100}),
      verificando qu\'e porcentaje de los verdaderos vecinos m\'as cercanos
      fueron recuperados exitosamente por \texttt{orange\_jumpsuit} frente a 
      FAISS. \footnote{
      Cabe aclarar que \textit{SiftSmall} usa $100$ vectores de busqueda,
      por ende brindando porcentajes de Recall\@n enteros.}
\end{itemize}


\begin{figure}[ht]
\begin{center}
\renewcommand{\arraystretch}{1.5}
{
    \rowcolors{2}{black!5}{black!10}
    \begin{tabularx}{0.8\textwidth} {
       >{\hsize=1\hsize\linewidth=\hsize}X 
       >{\hsize=1\hsize\linewidth=\hsize\centering\arraybackslash}X 
       >{\hsize=1\hsize\linewidth=\hsize\centering\arraybackslash}X }
    % \begin{tabular}{| c | m{5cm} | m{5cm} |} 
    \hline 

    \rowcolor{white}
     & \textbf{faiss} & \textbf{orange\_jumpsuit} \\

    \hline

    % \multicolumn{3}{l}{\hspace{10pt} Rendimiento (s)} \\
    Rendimiento (s) & & \\

    \hline 

    Train   & $5.5077$ & $7.9994$ \\
    Add     & $0.2758$ & $0.1743$ \\
    Search  & $0.1139$ & $0.1674$ \\

    \hline 

    % \multicolumn{3}{l}{M\'etricas (\%)} \\
    M\'etricas (\%) & & \\

    \hline 

    Recall@1    & $47\%$ & $46\%$ \\
    Recall@10   & $87\%$ & $83\%$ \\
    Recall@100  & $100\%$ & $100\%$ \\

    \hline 
    \end{tabularx}
    % \end{tabular}
}
\end{center}

\caption{\texttt{ComparingTest.ipynb} con el dataset \textit{SiftSmall}}
\label{fig:comparingtest_siftsmall}

\end{figure}
